\section{The exact generalized exceptional dimension}
\label{gen:section}

The generalized exceptional set allows arbitrary positive integer
multipliers in place of powers of two.  Following
\cite[Definition~1.5]{AL2014}, put
\[
 E_* = \{x\in\Sigma: Mx\in\Sigma\text{ for infinitely many }
                         M\in\mathbb Z_{>0},\ 3\nmid M\}.
\]
For a nonzero $x\in\Sigma$, every such multiplier is congruent to $1$
modulo $3$, as is seen from the first nonzero digit of $x$.
To specify the finite-visit convention, let
\[
 V_J=\{x\in\Sigma:\text{there are }1=M_1<\cdots<M_J,
       \ M_i\equiv1\pmod3,\ M_ix\in\Sigma\text{ for every }i\}.
\]
Write
\[
 d_*:=\frac{\log2}{2\log3},\qquad
 \Gamma_*:=\lim_{J\to\infty}\dimH V_J.
\]
The limit exists by monotonicity.  This is the generalized nesting
constant of \cite[(1.14)]{AL2014}; a shift in the indexing convention
for finite intersections does not change its limit.  It is distinct
from the dyadic constants $\Gamma_{\mathrm H},\Gamma_{\mathrm P}$.

\begin{theorem}[Exact generalized dimension]\label{gen:main}
One has
\[
 \boxed{\dimH E_* =\dimP E_* =\Gamma_*
       =\lim_{J\to\infty}\dimP V_J=d_*.}
\]
Moreover,
\begin{equation}\label{gen:rate-statement}
 \limsup_{J\to\infty}
       (\dimP V_J-d_*)\log\log J
       \le \frac{(\log2)^2}{\log3}.
\end{equation}
\end{theorem}

The lower bound $d_*$ was proved in
\cite[Theorems~2.4 and~6.3]{AL2014}, using a construction of
A.~Bolshakov.  We reproduce that construction below and do not claim
its novelty.  Our contribution here is the matching upper bound,
the equality with the generalized nesting constant, the packing
conclusions, and the quantitative estimate
\eqref{gen:rate-statement}.  The optimal rate of convergence is not
asserted.  Thus the result determines the dimension asked about in
\cite[Problem~1.6]{AL2014}, beyond the positive lower bound already
established there.

The proof uses a causal palette lemma specifically for pairs.  We
prove this finite lemma in full: no substitution of $r=2$ into a
container theorem requiring $r\ge3$ is made.  Throughout this section
all logarithms are natural.  Let $\Sigma_q$ be the image of $\Sigma$
in $\mathbb Z/3^q\mathbb Z$, represented in $[0,3^q)$, and write
$[z]_q$ for the canonical representative of $z\bmod3^q$.

\subsection{Capacity of adjacent translates}

\begin{lemma}[One-translate capacity]\label{gen:one-capacity}
For $A,B\subset\mathbb Z/3^q\mathbb Z$, if
$A+B\subset z+\Sigma_q$, then $|A||B|\le2^q$.
\end{lemma}
\begin{proof}
Empty sets cause no difficulty.  Translate $A$ to remove $z$ and
argue by induction on $q$, starting with the trivial case $q=0$.
The images of $A,B$ modulo $3$ have their sum in $\{0,1\}$.
Neither image can contain all three residues; if both contain two,
their sum is all of $\mathbb Z/3\mathbb Z$, since two subsets of
size two intersect after any translation.  Hence one image is a
singleton and the other has at most two elements.  There are at most
two nonempty pairs of least-digit sections.

For a pair of residues $a,b\in\{0,1,2\}$, write the corresponding
sections as $a+3A_a$ and $b+3B_b$.  Removing the least digit of their
sum leaves $A_a+B_b$ in a single translate of $\Sigma_{q-1}$;
the carry $\lfloor(a+b)/3\rfloor$ is fixed.  Thus induction bounds
$|A_a||B_b|$ by $2^{q-1}$.  Summing over the at most two pairs
proves the assertion.  All these operations are in the indicated
cyclic groups, so wraparound introduces no extra case.
\end{proof}

\begin{lemma}[Sharp capacity for two adjacent translates]
\label{gen:adjacent-capacity}
Let $q\ge1$, $A,B\subset\Sigma_q$ and $u\in\mathbb Z_3^\times$.
If
\[
 uA+B\subset (z+\Sigma_q)\cup(z+1+\Sigma_q)
                 \quad\text{in }\mathbb Z/3^q\mathbb Z,
\]
then $|A||B|\le2^{q+1}$.  This bound is sharp.
\end{lemma}
\begin{proof}
Replace $uA$ by $uA-z$.  Each of $uA-z$ and $B$ has at most two
residues modulo $3$.  The exact identity
\[
 \Sigma_q\cup(1+\Sigma_q)=\{0,1,2\}+3\Sigma_{q-1}
\]
shows that, for each of the at most four pairs of least-digit
sections, the sum of the upper sections lies in a single translate
of $\Sigma_{q-1}$.  Lemma~\ref{gen:one-capacity} bounds each product
of section sizes by $2^{q-1}$.  Multiplication by $u$ and translation
preserve cardinalities, so their sum gives $2^{q+1}$.
Equality is attained by $u=1$, $z=0$, $A=\{0,1\}$ and $B=\Sigma_q$.
Adjacency of the two translates is essential to this argument.
\end{proof}

\subsection{A causal palette lemma for pairs}

The following statement permits overlapping windows.  In particular,
it does not assume that the letters associated with different
positions are independent.

\begin{lemma}[Causal pair palettes]\label{gen:palettes}
Let $W$ be a nonempty finite set of input words of length $L$.
For $0\le j<N\le L$, assign to each $x\in W$ a letter
$Z_j(x)\in\Omega$, where $|\Omega|=Q$.  Let a fixed simple directed
graph have edges $(j,k)$ with $j<k$.  For each edge let
$\mathcal R_{jk}(x_{<j})\subset\Omega^2$ depend only on the input
prefix preceding position $j$, and suppose that
\[
 (Z_j(x),Z_k(x))\in\mathcal R_{jk}(x_{<j})
             \qquad(x\in W).
\]
For an integer $T\ge2Q$, put $\theta=Q/T$.  Each $x\in W$ admits
a deterministic fingerprint $F(x)$, consisting of at most
$\lfloor\theta N\rfloor$ pairs $(j,Z_j(x))$, and nonempty palettes
$A_j(x)\subset\Omega$ with the following properties:
\begin{enumerate}[label=(\roman*)]
\item $Z_j(x)\in A_j(x)$, and $A_j(x)$ is determined by
      $(x_{<j},F(x))$;
\item $A_j(x)\times A_k(x)\subset\mathcal R_{jk}(x_{<j})$
      except on at most $QTN$ edges;
\item for any probability distribution on $W$,
\[
 H(F)\le N\{h(\theta)+\theta\log Q\},\qquad
 h(t)=-t\log t-(1-t)\log(1-t).
\]
\end{enumerate}
Here $H$ denotes Shannon entropy with natural logarithms.
No locality assumption on the maps $Z_j$ is required.
\end{lemma}
\begin{proof}
Initially give each position the mutable palette $S_j=\Omega$.
Process source positions in increasing order.  At source $j$, call
$a\in S_j$ heavy if at least $T$ outgoing edges $(j,k)$ admit some
$b\in S_k$ with $(a,b)\notin\mathcal R_{jk}(x_{<j})$.

If the actual letter $z_j=Z_j(x)$ is heavy, record $(j,z_j)$ in
$F$, set $A_j=\{z_j\}$, and remove from every future $S_k$ all
letters forbidden with source $z_j$ on the edge $(j,k)$, if that edge
exists.  These removals never delete the actual future letters.
At least $T$ distinct future letters are removed: the witnessing
edges have distinct targets because the graph is simple.
If $z_j$ is light, let $A_j$ be the set of all light letters in
the current $S_j$, and make no removals.  Once assigned, $A_j$ is
fixed; only palettes at unprocessed positions can shrink.

There are at most $QN$ possible removals in total, so
$|F|T\le QN$.  On every edge from a heavy source, the full product
of the final palettes is allowed.  At a light source, each of its
at most $Q$ assigned letters has fewer than $T$ problematic edges.
Subsequent shrinking cannot create new forbidden pairs.  This gives
at most $QT$ bad edges per light source and at most $QTN$ overall.

For replay at position $j$, use the known prefix $x_{<j}$ and the
whole fingerprint $F$.  At every earlier source $i\le j$, all
required relations are known from the corresponding prefix.  Source
$i$ is recorded exactly when its actual letter was heavy, in which
case $F$ supplies that letter.  If it is not recorded, the algorithm
needs only the computable set of light letters, not the unknown
actual letter.  Thus replay recovers $A_j$ from the stipulated data.

Finally the number of possible fingerprints is at most
\[
 \sum_{m\le\lfloor\theta N\rfloor}\binom Nm Q^m
 \le\exp\bigl(Nh(\theta)+\theta N\log Q\bigr).
\]
For completeness, when $m\le\theta N$ and $\theta\le1/2$,
$\theta^m(1-\theta)^{N-m}\ge e^{-Nh(\theta)}$.
Summing the corresponding terms in the binomial identity proves
$\sum_{m\le\theta N}\binom Nm\le e^{Nh(\theta)}$.
The logarithm of the number of fingerprints bounds their entropy.
\end{proof}

\subsection{Uniform affine counting from two windows}

Fix distinct positive integers $E_1,\ldots,E_K$ and units
$u_i\in\mathbb Z_3^\times$, and put
\[
 c_i=1+3^{E_i}u_i,\qquad R=\max_i E_i.
\]
The shifts in the next lemma are arbitrary $3$-adic integers.  Their
uniform treatment will also permit localization to every cylinder.

\begin{lemma}[Finite affine bound]\label{gen:affine-count}
Let $q\ge1$, $Q=2^q$, $T\ge2Q$ an integer, and $\theta=Q/T$.
Define
\begin{equation}\label{gen:sigma}
 \sigma=d_*+\frac{d_*}{q}
       +\frac{h(\theta)}{\log3}
       +\theta q\frac{\log2}{\log3}
       +\frac{QT}{K}\frac{\log2}{\log3}.
\end{equation}
For $\beta_1,\ldots,\beta_K\in\mathbb Z_3$, let
\[
 W_L(\beta)=\{X\in\Sigma_L:
                   [c_iX+\beta_i]_L\in\Sigma_L\ (1\le i\le K)\}.
\]
For every $L\ge q$, with $N=L-q+1$, a nonempty $W_L(\beta)$ satisfies
\begin{equation}\label{gen:finite-log}
 \log |W_L(\beta)|
       \le N\sigma\log3+(R+q-1)\log2.
\end{equation}
In particular, whether or not it is empty,
\begin{equation}\label{gen:finite-count}
 |W_L(\beta)|\le2^{R+q-1}3^{\sigma L}.
\end{equation}
These bounds are uniform in all the shifts $\beta_i$.
\end{lemma}
\begin{proof}
Assume $W_L(\beta)$ is nonempty.  For $0\le j<N$, use the window
letter
\[
 Z_j(X)=B_{j,q}(X):=\left\lfloor X/3^j\right\rfloor\bmod3^q
                    \in\Sigma_q.
\]
For each $i$ include the edge $(j,v)=(j,j+E_i)$ whenever $v<N$.
The distinctness of the $E_i$ makes this a simple graph.  For this
edge set
\[
 P=[X]_j,\quad A=\beta_i+3^{E_i}u_iP,\quad
 a_v=[A]_v,\quad D=\frac{A-a_v}{3^v}\in\mathbb Z_3.
\]
Here $D$ depends only on the fixed data and the prefix before $j$.
An exact decomposition modulo $3^{v+q}$ gives
\[
 c_iX+\beta_i\equiv [X]_v+a_v+
          3^v\bigl(B_{v,q}(X)+u_iB_{j,q}(X)+D\bigr).
\]
Both $[X]_v$ and $a_v$ are ordinary integers in $[0,3^v)$.
Their carry is therefore $e\in\{0,1\}$, and admissibility implies
\begin{equation}\label{gen:two-window-relation}
 u_iB_{j,q}(X)+B_{v,q}(X)+D+e\in\Sigma_q\pmod{3^q}
                 \quad\text{for some }e\in\{0,1\}.
\end{equation}
This argument is valid for arbitrary $3$-adic $u_i$ and $\beta_i$:
only their canonical lower residues enter the ordinary carry.

Apply Lemma~\ref{gen:palettes} with $\Omega=\Sigma_q$ and relations
\eqref{gen:two-window-relation}.  On a good edge the full product
of palettes lies in this relation.  Its image under $(a,b)\mapsto
u_i a+b$ lies in two adjacent translates of $\Sigma_q$, so
Lemma~\ref{gen:adjacent-capacity} gives
\[
 w_j+w_v\le(q+1)\log2,
       \qquad w_j:=\log|A_j|\le q\log2.
\]
On a bad edge use $w_j+w_v\le2q\log2$.  The graph has at most
$KN$ edges, at most $QTN$ of them bad.  Compared with two incidences
at each position for each shift, shift $E_i$ loses at most
$2\min(E_i,N)\le2R$ incidences.  Consequently, for each input word,
\begin{align}
 2K\sum_{j=0}^{N-1}w_j
 &\le KN(q+1)\log2+2q\log2\,QTN+2KRq\log2,\notag\\
 \sum_{j=0}^{N-1}w_j
 &\le \frac{N(q+1)}2\log2
             +q\log2\left(\frac{QTN}{K}+R\right).
 \label{gen:weights}
\end{align}
Using $KN$ in the first term only enlarges the bound, including
when some shifts exceed the window range.

Now choose $X$ uniformly in $W_L(\beta)$ and write
$h_t=H(X_t\mid X_{<t},F)$.  Replay in
Lemma~\ref{gen:palettes} ensures that, conditional on $(X_{<j},F)$,
the support of $B_{j,q}(X)$ lies in the corresponding palette.
The entropy chain rule therefore gives
\[
 \sum_{t=j}^{j+q-1}h_t
 =H(B_{j,q}(X)\mid X_{<j},F)\le\mathbb E w_j.
\]
This identity, rather than independence of overlapping windows, is
the reason their simultaneous use is legitimate.  Each digit is
covered at most $q$ times; the total deficit from coverage $q$ at
every digit is $q(q-1)$.  Since $h_t\le\log2$,
\[
 qH(X\mid F)\le\sum_j\mathbb E w_j+q(q-1)\log2.
\]
The fingerprint is a deterministic function of $X$, whence
$H(X)=H(F)+H(X\mid F)$.  Substituting
\eqref{gen:weights} and the fingerprint entropy bound gives
\eqref{gen:finite-log}, with precisely \eqref{gen:sigma}.
Finally $N\le L$ and $\sigma\ge0$ give
\eqref{gen:finite-count}; the empty case is immediate.
\end{proof}

\subsection{Finite extraction and uniform local bounds}

\begin{lemma}[Distinct depths from a finite unit set]
\label{gen:unit-tree}
A set of more than $2\cdot3^{K-1}$ distinct elements of
$\mathbb Z_3^\times$ contains an anchor $\lambda_0$ and
$\lambda_1,\ldots,\lambda_K$ such that
$v_3(\lambda_i-\lambda_0)$ are distinct positive integers.
\end{lemma}
\begin{proof}
One nonzero residue class modulo $3$ contains more than $3^{K-1}$
elements.  Form their tree of successive residues, stopping at a
finite depth at which all elements are separated; such a depth
exists because the set is finite and its elements are distinct.
Suppress vertices with a single child, retaining the original
depth at every branching vertex.  The resulting tree has at most
three children per vertex.  If every root-to-leaf path had at most
$K-1$ branchings, induction would give at most $3^{K-1}$ leaves.
Choose a path with at least $K$ branchings and let its terminal
leaf be $\lambda_0$.  At each of $K$ branchings choose a leaf in
a different child subtree.  The first differing positions of these
leaves from $\lambda_0$ are exactly the distinct branching depths.
They are positive because all leaves agree modulo $3$.
\end{proof}

We use the usual Assouad dimension $\dim_{\mathrm A}$: an exponent
$a$ is admissible if every intersection of a set with a ball of
radius $\rho$ can be covered by at most $C(\rho/r)^a$ balls of
radius $r$, with one $C$ for all centers and $0<r<\rho\le1$.

\begin{theorem}[Uniform finite intersections]\label{gen:finite-uniform}
For every $\varepsilon>0$ there is an integer $J(\varepsilon)$ such
that, for any $J\ge J(\varepsilon)$ distinct units
$\lambda_1,\ldots,\lambda_J\in\mathbb Z_3^\times$,
\[
 \dim_{\mathrm A}\bigcap_{i=1}^J\lambda_i^{-1}\Sigma
                \le d_*+\varepsilon.
\]
In particular the same bound holds for upper box, packing and
Hausdorff dimensions.  The covering constant may depend on the
chosen multipliers; the exponent and sufficient cardinality do not.
\end{theorem}
\begin{proof}
Choose $q$, then $T$, then $K$ so that the three nonconstant terms
after $d_*$ in \eqref{gen:sigma} total less than $\varepsilon$:
first $d_*/q<\varepsilon/3$, next
\[
 \frac{h(Q/T)}{\log3}+\frac QT q\frac{\log2}{\log3}
              <\frac\varepsilon3,
\]
and finally $(QT/K)(\log2/\log3)<\varepsilon/3$.
Take $J(\varepsilon)>2\cdot3^{K-1}$.  By
Lemma~\ref{gen:unit-tree}, choose an anchor and $K$ companions.
Multiplication by the anchor is an isometry; the image of the full
intersection is contained in
\[
 G=\{y\in\Sigma:c_i y\in\Sigma\ (1\le i\le K)\},
       \qquad c_i=1+3^{E_i}u_i,
\]
with the distinct positive depths required in
Lemma~\ref{gen:affine-count}.

Consider any cylinder $P+3^t\mathbb Z_3$ meeting $G$, with
$0\le P<3^t$.  Its lower input and output digits are already
admissible.  Writing $y=P+3^tz$ yields the exact tail conditions
\[
 z\in\Sigma,\qquad c_i z+\beta_i\in\Sigma,
       \quad\beta_i=\frac{c_iP-[c_iP]_t}{3^t}\in\mathbb Z_3.
\]
Thus the number of depth-$t+L$ cylinders meeting $G$ inside it is
at most $2^{R+q-1}3^{\sigma L}$ for $L\ge q$, uniformly in $P,t$.
For $0\le L<q$ the elementary bound $2^L\le2^{R+q-1}$ suffices.
This proves the local covering estimate on every pair of discrete
$3$-adic scales, with one constant.  Rounding arbitrary radii to
neighboring powers of $3$ changes the constant only by a factor
depending on $\sigma$.  Hence $\dim_{\mathrm A}G\le\sigma<d_*+
\varepsilon$.  Monotonicity and the anchor isometry finish the proof.
\end{proof}

\subsection{The exact value and its convergence rate}

\begin{proof}[Proof of Theorem~\ref{gen:main}]
For a sufficiently large fixed $J$, the set $V_J$ is a countable
union, over integer multiplier tuples, of the compact intersections
in Theorem~\ref{gen:finite-uniform}.  Therefore
$\dimP V_J\le d_*+\varepsilon$ by countable stability of packing
dimension.  No countable stability of upper box or Assouad dimension
is used.  Since $E_*\subset V_J$ for every $J$, this proves
\[
 \dimP E_*\le d_*,\qquad
 \Gamma_*\le d_*,\qquad
 \lim_{J\to\infty}\dimP V_J\le d_*.
\]
The inclusion also covers zero; at a nonzero point the permissible
unit multipliers are $1\pmod3$, and the multiplier $1$ is available
because the point lies in $\Sigma$.

For the known lower bound, consider
\[
 U=\left\{\sum_{j\ge0}\epsilon_j3^{2j}:\epsilon_j\in\{0,1\}\right\}.
\]
For every $k\ge0$, the supports of $x\in U$ and $3^{2k+1}x$
have opposite parities, so $(1+3^{2k+1})x\in\Sigma$ without carry.
Thus $U\subset E_*\cap V_J$ for every $J$.  This is the
Bolshakov construction used in \cite[Theorem~6.3]{AL2014}.
There are $2^{\lceil L/2\rceil}$ depth-$L$ cylinders of $U$.
The fair Bernoulli measure on the even digits assigns each such
cylinder mass $2^{-\lceil L/2\rceil}\le3^{-d_*L}$.
The cylinder cover and the mass distribution principle give
$\dimH U=\dimP U=d_*$.  This proves the asserted exact values.

For the rate, take, for all sufficiently large $J$,
\[
 K=\left\lfloor\log_3\frac{J-1}{2}\right\rfloor+1,
 \qquad
 q=\left\lfloor\frac{\log K-5\log\log K}{\log4}\right\rfloor,
 \qquad Q=2^q,\quad T=Qq^3.
\]
Then $q\ge2$, $\theta=q^{-3}\le1/2$, and
$2\cdot3^{K-1}\le J-1<J$, so finite extraction is available for
every tuple in the definition of $V_J$.  Moreover,
\[
 h(\theta)=O(q^{-3}\log q),\qquad
 \theta q=q^{-2},\qquad
 \frac{QT}{K}=\frac{4^q q^3}{K}
       \le\frac{q^3}{(\log K)^5}=O(q^{-2}).
\]
Consequently \eqref{gen:sigma} gives
$\sigma=d_*+d_*/q+O(q^{-2})$.  The same countable-union argument
using \eqref{gen:finite-count} implies $\dimP V_J\le\sigma$.
Finally $q\sim\log K/\log4$, $\log K\sim\log\log J$, and
$d_*\log4=(\log2)^2/\log3$.  This proves
\eqref{gen:rate-statement}.  The constants involving $R$ disappear
in the dimension estimate for each fixed tuple before the countable
union is taken.
\end{proof}

\subsection{A sharp local bound for fixed infinite families}

\begin{theorem}[Fixed infinite intersections]\label{gen:assouad}
For every infinite set of distinct units
$\Lambda\subset\mathbb Z_3^\times$,
\[
 F_\Lambda:=\bigcap_{\lambda\in\Lambda}\lambda^{-1}\Sigma
       \quad\text{satisfies}\quad
 \dim_{\mathrm A}F_\Lambda\le d_*.
\]
The bound is sharp.  For
$\Lambda_{\rm odd}=\{1\}\cup\{1+3^{2k+1}:k\ge0\}$,
\[
 F_{\Lambda_{\rm odd}}=U\cup3U,
       \qquad\dim_{\mathrm A}F_{\Lambda_{\rm odd}}=d_*.
\]
\end{theorem}
\begin{proof}
For each $\varepsilon>0$, choose any $J(\varepsilon)$ distinct
members of $\Lambda$.  The full intersection is contained in
this finite intersection, so Theorem~\ref{gen:finite-uniform} and
monotonicity give $\dim_{\mathrm A}F_\Lambda\le d_*+\varepsilon$.
One fixed subfamily for each exponent suffices; no uniform covering
constant over all finite subfamilies, nor countability of $\Lambda$,
is required.

In the displayed example, adding $x$ and $3^h x$ never creates a
carry, since each sum of digits is at most two.  Their sum omits
the digit two exactly when no two nonzero digits of $x$ are distance
$h$ apart.  Requiring this for every positive odd $h$ forces all
nonzero digits to have the same parity.  This is precisely
$x\in U\cup3U$.  In each cylinder of $U$, the number of descendants
$L$ levels below is at most $2^{\lceil L/2\rceil}$; hence
$\dim_{\mathrm A}U\le d_*$.  The reverse inequality follows
from its Hausdorff dimension, and the stated sharpness follows.
\end{proof}

\begin{remark}[The family must be fixed]\label{gen:scope}
Theorem~\ref{gen:assouad} does not imply
$\dim_{\mathrm A}E_*\le d_*$.  Every finite $0,1$-digit integer
$x$ belongs to $E_*$: for arbitrarily large $h$, the finite supports
of $x$ and $3^hx$ are disjoint, so $(1+3^h)x\in\Sigma$.
These integers are dense in $\Sigma$.  Thus the upper box dimension
of $E_*$ is $\log_3 2$, and monotonicity together with
$\overline{\dim}_{\mathrm B}E_*\le\dim_{\mathrm A}E_*$ gives
\[
 \overline{\dim}_{\mathrm B}E_*
      =\dim_{\mathrm A}E_*=\log_3 2.
\]
For $E_*$ the multiplier family can depend on the point; for
$F_\Lambda$ the same family works for every point.  In particular
$1\in E_*$, and none of these generalized results excludes
$1$ from the dyadic exceptional set $E_2$.
\end{remark}
